% from ega0-21-fr.tex
\oldpage[0\textsubscript{IV-1}]{166}
ce dernier se factorise donc en $A\xrightarrow{u'}B\to B/\mathfrak n^2$ et
les hypothèses sur $\operatorname{gr}_0(u'_1)$ et
$\operatorname{gr}_1(u'_1)$ entraînent
$\operatorname{gr}_0(u')=\operatorname{gr}_0(u)=u_0$ et
$\operatorname{gr}_1(u')=\operatorname{gr}_1(u)$, donc l'image par $u'$ d'une
uniformisante de $A$ est une uniformisante de $B$, et l'on a par suite
$Bu'(\mathfrak m)=\mathfrak n$. Il suffit alors, pour obtenir un automorphisme
$w_1$ répondant à la question, de prendre l'automorphisme déduit par passage
aux quotients de l'automorphisme $w$ de $B$ fourni par l'application de $b')$
à l'homomorphisme $u'$.
\end{proof}
\begin{remarks}[21.5.3]
\label{0.21.5.3-fr}
\begin{enumerate}
  \item[(i)] Les propriétés différentielles des corps permettent de résoudre la
  question de l'unicité du corps de représentants dans un anneau local
  noethérien complet $C$ (19.8.7, (ii)). Soient $\mathfrak J$ l'idéal maximal
  de $C$, $K=C/\mathfrak J$ le corps résiduel de $C$~; on peut se borner au
  cas $\mathfrak J\ne0$. Supposons qu'il existe un homomorphisme $u:K\to C$
  qui, composé avec l'augmentation $C\to K$, donne l'identité~; alors, pour
  que cet homomorphisme soit unique, il faut et il suffit que
  $\Omega^1_K=0$. En effet, la condition $\Omega^1_K=0$ entraîne que $K$ est
  formellement non ramifié sur son corps premier (20.7.4)~; s'il existait un
  second homomorphisme $v\ne u$ répondant à la question, il y aurait un plus
  grand entier $n$ tel que $\mathfrak J^n$ contienne l'ensemble des
  $u(x)-v(x)$ pour $x\in K$~; par passage au quotient, $u$ et $v$ donneraient
  deux homomorphismes distincts $u_1,v_1$ de $K$ dans
  $C/\mathfrak J^{n+1}$, dont les composés avec
  $C/\mathfrak J^{n+1}\to C/\mathfrak J^n$ seraient égaux, ce qui contredit
  la définition (19.10.2). Inversement, supposons que $\Omega^1_K\ne0$~; il
  existe alors une dérivation $D\ne0$ de $K$ dans
  $\mathfrak J/\mathfrak J^2$ (20.4.8), donc un homomorphisme
  $v_1:K\to C/\mathfrak J^2$ tel que, si $u_1:K\to C/\mathfrak J^2$ est
  obtenu par passage au quotient à partir de $u$, on ait $v_1=u_1+D$
  (20.1.1). Si $k$ est le corps premier de $K$, $C$ est une $k$-algèbre et
  $K$ une $k$-algèbre formellement lisse (19.6.1), et les restrictions de
  $u_1$ et $v_1$ à $k$ coïncident, donc $v_1$ se factorise en
  $K\xrightarrow{v}C\to C/\mathfrak J^2$ et l'on a $v\ne u$.
  Rappelons ((20.6.20) et (21.4.4)) que la condition $\Omega^1_K=0$ signifie
  que $K$ est parfait s'il est de caractéristique $\ne0$, et extension
  algébrique de $\mathbf Q$ s'il est de caractéristique $0$.
  \item[(ii)] De la même manière, soient $W$ un anneau de Cohen, $C$ un anneau
  local noethérien complet, soit $\mathfrak J\ne0$ un idéal de $C$ contenu dans
  l'idéal maximal~; pour que la factorisation $W\to C\to C/\mathfrak J$
  dans (19.8.6, (i)) soit unique, il faut et il suffit que le corps résiduel
  $K$ de l'anneau de Cohen $W$ soit tel que $\Omega^1_K=0$. En effet, si
  $\Omega^1_K\ne0$ et si $\mathfrak m$ désigne l'idéal maximal de $W$, il
  suffit de composer une dérivation $D\ne0$ de $K$ dans
  $\mathfrak J/\mathfrak m\mathfrak J$ avec l'augmentation $W\to K$ pour
  obtenir une dérivation $D'\ne0$ de $W$ dans
  $\mathfrak J/\mathfrak m\mathfrak J$ et l'on termine le raisonnement comme
  dans (i) en formant à l'aide de $D'$ un homomorphisme
  $v_1:W\to C/\mathfrak m\mathfrak J$ distinct de l'homomorphisme
  $u_1:W\to C/\mathfrak m\mathfrak J$ déduit de $u$ par passage au quotient~;
  on achève le raisonnement en invoquant cette fois (19.8.6, (i)). Si au
  contraire $\Omega^1_K=0$, l'unicité de $v$ résulte déjà de (i) lorsque
  $W=K$ est un corps de caractéristique $0$. Dans le cas contraire, on a
  $\widehat{\Omega}^1_W=0$~; en effet, on a alors $\mathfrak m=pW$, $p$ étant
  la caractéristique de $K$ (19.8.5), et l'homomorphisme canonique
  $\mathfrak m/\mathfrak m^2\to\Omega^1_W/\mathfrak m\Omega^1_W$ (20.5.11.2)
  est par suite nul. La suite exacte (20.5.12.1) appliquée à $W$ et à
  $K=W/\mathfrak m$ entraîne alors que $\Omega^1_W=\mathfrak m\Omega^1_W$,
  d'où notre assertion. Mais alors (20.7.4) $W$ est formellement non ramifié
  (pour sa topologie adique) sur $\mathbf Z$, et l'unicité de $v$ se prouve
  comme dans (i).
\end{enumerate}
\end{remarks}

